\section{可換三生成部分上のOreチャートと第三階数}
\label{sec:polynomial-ore}
低階数残余を調べるため、三生成部分が可換な場合を一括して扱う。
以下は既知のOre構成を用いる共通の認証チャートであり、全対象の被覆ではない。

$B=k[x_0,x_1,x_2]$とし、任意の二次式
\[
 f_i=\sum_{0\leq a\leq b\leq2}c_{i,ab}x_ax_b\quad(i=0,1,2)
\]
を取る。$\delta(x_i)=f_i$で定める導分に対し
\[
 A_f=B[X;\delta],\qquad Xx_i-x_iX=f_i
\]
と置く。係数空間は$\mathbb A^{18}$である。

\begin{theorem}[全ファイバー正則な多項式Oreチャート]\label{thm:polynomial-ore-regular}
全ての係数で$A_f$は標準14641型quadratic Koszul AS正則であり、
$\check A_f$は1周期的である。基底は$x_0^{a_0}x_1^{a_1}x_2^{a_2}X^m$である。
これは18次元の同型類像や全モジュライ成分を意味しない。
\end{theorem}
\begin{proof}
任意の三つの$f_i$は多項式環の導分を定める。
$x_0<x_1<x_2<X$の書換えで$Xx_i$を$x_iX+f_i$に直すと、
$X$の個数を第一重み、反転数を第二順序として停止する。
三次重なりは$Xx_jx_i$の三つと$x_2x_1x_0$の一つで、
$\delta(x_jx_i-x_ix_j)=0$により全て解消する。
従って二次Gr\"obner基底であり、主張の基底とKoszul性を得る。

双対生成元を$e_0,e_1,e_2,z$とすると、二次双対は
\begin{align*}
 z^2&=0,& ze_i+e_iz&=0,\\
 e_i^2&=\sum_a c_{a,ii}e_az,&
 e_ie_j+e_je_i&=\sum_a c_{a,ij}e_az\quad(i<j)
\end{align*}
である。$e_a z$は中心的であり、二次補正をもう一度使うと$z^2$が現れる。
$e_i$に重み1、$z$に重み0を与えた書換えは、平方自由の昇順語へ停止する。
20個の三次重なりは恒等的に解消し、随伴次数代数は四生成外積代数となる。
この20恒等式は\path{computations/verify_polynomial_ore.py}で18変数のまま検査する。
従って双対は次数元数14641を持ち、最上積$e_0e_1e_2z$の係数をtraceとする相補積対合は
外積代数の対合のフィルター付き持上げとして完全である。
Koszul性とこの有限双対Frobenius性から、長さ4の線型解消と左右AS条件を得る。
\end{proof}

$w_f$をこの最上traceの引き戻しとし、$z$を第3生成元と読む。
最初の四一次空間の基底を保って標準4周期表示へ直せるので、
この$w_f$の第三階数は内在的な$\rho$に一致する。
少なくとも二つの$z$を含む四次語のtraceは零であり、ちょうど一つの場合は外積の符号である。
$z$を含まない四次語では二次補正が一度だけ寄与するので、そのtraceは$c$に線型である。
$\mathcal P=(00,01,02,11,12,22)$の順に、$U=\langle x_0,x_1,x_2\rangle$の
対称二次テンソルを並べる。各順序対$ac$を$e_a\otimes e_c$と同一視し、
非対角対には両順序を一度ずつ足す規約で
\[
 Q_{ac,bd}(f)=\sum_{(a',c')\in\operatorname{ord}(ac)}
              \sum_{(b',d')\in\operatorname{ord}(bd)}
              (w_f)_{a'b'c'd'}
\]
を定める。これは$6\times6$行列である。

\begin{theorem}[第三階数の交代行列判定]\label{thm:polynomial-ore-rank}
$Q(f)$は交代行列で、$f\mapsto Q(f)$は18次元係数空間から交代$6\times6$行列空間への全射線型写像である。
その核はちょうど$f_i=x_i\ell$（$\ell$は任意の一次式）となる3次元部分空間であり、
全パラメータで
\[
 \rho(\check A_f)=6+\rk Q(f)
\]
が成立する。従ってこのチャートの第三階数は正確に$6,8,10,12$である。
$\rho\leq6$の全ての対象は可換多項式代数のindexed類に属し、
このチャートは\texttt{LOW-UNCLASSIFIED}に対象を与えない。
\end{theorem}
\begin{proof}
第三の二対二行列で、両因子が$U$に属する9成分と、一方が第3生成元である6成分を分けると、
零行・零列を除いて
\[
 M_f=\begin{pmatrix}D(f)&B_0\\ C_0&0\end{pmatrix}
\]
となる。$B_0,C_0$は外積の符号だけで決まり、双方の階数は3である。
$\operatorname{im}B_0=\Lambda^2U$、$\ker C_0=\Sym^2U$なので、
$D(f)$の$\ker C_0\to\operatorname{coker}B_0$への誘導行列が、定義した$Q(f)$である。
従ってブロック消去により$\rk M_f=3+3+\rk Q(f)$となる。

$\mathcal P$を$0,\ldots,5$と番号付けすると、上三角成分は次の通りである。
\[
\begin{array}{c|l@{\qquad}c|l}
 ij&Q_{ij}&ij&Q_{ij}\\\hline
 01&2c_{2,00}&12&2c_{0,00}-c_{1,01}-c_{2,02}\\
 02&-2c_{1,00}&13&2c_{2,11}\\
 03&c_{2,01}&14&c_{0,01}-2c_{1,11}+c_{2,12}\\
 04&-c_{1,01}+c_{2,02}&15&c_{0,02}-c_{1,12}\\
 05&-c_{1,02}&23&-c_{0,01}+c_{2,12}\\
 24&-c_{0,02}-c_{1,12}+2c_{2,22}&25&-2c_{1,22}\\
 34&2c_{0,11}&35&c_{0,12}\\
 45&2c_{0,22}&&
\end{array}
\]
対角成分は零、下三角成分は上三角の負である。
双対の書換えからこの表を直接計算でき、再現コードは四次の全256語で同じ表を検査する。
$Q=0$は、まず$i\notin\{a,b\}$の九係数が零であることを強制し、
残る九係数には
\[
 c_{0,00}=c_{1,01}=c_{2,02},\qquad
 c_{0,01}=c_{1,11}=c_{2,12},\qquad
 c_{0,02}=c_{1,12}=c_{2,22}
\]
だけを課す。従って核は$f_i=x_i\ell$で、次元は3である。
線型写像の階数は15となり、交代行列空間全体への全射性が従う。
交代行列の階数は$0,2,4,6$なので、それぞれの非空性も含めて階数の主張を得る。

最後に$f_i=x_i\ell$なら、位置$j$の一次写像
$x_i\mapsto x_i$, $X\mapsto X-j\ell$を取る。
隣接位置の積に適用すると
\[
 (X-j\ell)x_i-x_i(X-(j+1)\ell)-\ell x_i=0
\]
が可換多項式環内で成立するため、これは全整数で定義されたindexed同型である。
逆に$\rho\leq6$なら階数式から$Q=0$となるので、この場合で尽きる。
\end{proof}
\begin{remark}
$Q$の階数で与える条件はこの係数空間のPfaffian階数条件であり、
indexed同型の完全不変量でも、全モジュライの成分方程式でもない。
Ore構成・PBW・有限双対の手法自体について新規性を主張しない。
\end{remark}
